% =====================================================================
\chapter{Actions dangereuses et confirmation humaine}
\label{ch:actions}

Ce chapitre recense \emph{toutes} les fonctions du dépôt capables de modifier un
système, de supprimer quelque chose, de démarrer un service, de toucher à un
volume, à un snapshot, à une sauvegarde ou à une configuration --- puis détaille
le mécanisme de confirmation.

% =====================================================================
\section{Inventaire des fonctions à effet de bord}

\subsection{Groupe A --- Modification de l'infrastructure (VM)}

\begin{longtable}{@{}p{3.6cm}p{3.5cm}p{7.3cm}@{}}
\toprule
\textbf{Fonction} & \textbf{Fichier} & \textbf{Ce qu'elle fait réellement}\\
\midrule
\endhead
\fn{run\_backup(target)} & \file{tools/backup.py} & \texttt{systemctl start <service>} --- démarre un job de sauvegarde existant\\
\fn{create\_snapshot(volume\_id, volume\_group, size\_gb)} & \file{tools/snapshot.py} & \texttt{sudo lvcreate -{}-snapshot -{}-size <n>G -{}-name <nom> /dev/<vg>/<lv>}\\
\fn{restore\_from\_backup(target, backup\_id, confirm)} & \file{tools/restore.py} & \texttt{sudo lvconvert -{}-merge <snapshot>} --- \textbf{écrase le volume}\\
\fn{repair\_backup\_repository(target, confirm)} & \file{tools/repair.py} & Écrit un drop-in systemd et recharge le démon\\
\fn{\_write\_drop\_in(service, body)} & \file{tools/repair.py} & \texttt{sudo -n install -d}, \texttt{sudo -n tee}, \texttt{systemctl daemon-reload}\\
\fn{ssh.run(command, allow\_failure)} & \file{mcp/storage/ssh.py} & \textbf{Exécute une commande arbitraire} sur la VM\\
\fn{ssh.run\_many(commands)} & \file{mcp/storage/ssh.py} & Idem, plusieurs commandes sur une session\\
\bottomrule
\end{longtable}

\begin{keybox}
\textbf{Point de sécurité important.} \fn{ssh.run} est la seule porte vers la
VM, et \textbf{aucune commande n'est composée à partir d'une chaîne fournie par
le modèle}. Le modèle fournit des \emph{arguments} (un nom d'unit, un nom de
volume), jamais un verbe. Les commandes sont écrites en dur dans les outils, et
chaque argument passe par \fn{quote} (\texttt{shlex.quote}) ou par une
conversion numérique.
\end{keybox}

\subsection{Groupe B --- Suppressions applicatives}

\begin{longtable}{@{}p{3.9cm}p{3.4cm}p{7.1cm}@{}}
\toprule
\textbf{Fonction} & \textbf{Fichier} & \textbf{Ce qu'elle supprime}\\
\midrule
\endhead
\fn{delete\_conversation} & \file{services/chat\_service.py} & La conversation et ses messages (\texttt{ON DELETE CASCADE})\\
\fn{delete\_document} & \file{services/knowledge\_service.py} & Les chunks dans Chroma, \emph{puis} le fichier sur disque, \emph{puis} la ligne en base\\
\fn{delete\_request} & \file{services/request\_service.py} & La demande, ses messages, son journal, ses pièces jointes, et les binaires sur disque\\
\fn{delete\_attachment} & \file{services/request\_service.py} & Une pièce jointe et son binaire\\
\fn{purge\_old} & \file{services/metrics\_service.py} & Les \tbl{disk\_metrics} de plus de 31 jours\\
\bottomrule
\end{longtable}

\begin{infobox}
\textbf{Aucune de ces suppressions ne touche l'infrastructure.} Elles restent
dans le périmètre applicatif. Elles ne passent donc pas par le mécanisme
d'\cls{ActionRequest} --- mais elles sont toutes protégées par
\fn{require\_role} et, pour les demandes et conversations, par un contrôle de
propriété.
\end{infobox}

\subsection{Groupe C --- Modifications d'état applicatif}

\fn{confirm\_and\_execute\_action}, \fn{reject\_action\_request},
\fn{resolve\_ticket}, \fn{update\_ticket\_status}, \fn{update\_status}
(demandes), \fn{approve\_request} / \fn{reject\_request} (cross-domain).
Toutes écrivent une trace : note d'incident ou ligne \tbl{request\_events}.

% =====================================================================
\section{Fiche détaillée de chaque action d'infrastructure}

% ---------------------------------------------------
\subsection{\fn{run\_backup(target)}}

\begin{longtable}{@{}p{3.2cm}p{11.6cm}@{}}
\toprule
\endhead
\textbf{Où} & \file{app/mcp/storage/tools/backup.py}\\
\textbf{Qui peut l'appeler} & Le modèle, mais l'appel est \textbf{intercepté} : \texttt{run\_backup} $\in$ \texttt{REMEDIATION\_TOOLS}. Exécution réelle uniquement via \route{POST /storage/actions/\{id\}/confirm}\\
\textbf{Paramètres} & \texttt{target: str} --- nom du timer (\texttt{backup-app-critical.timer}) ou du service\\
\textbf{Effet réel} & Convertit \texttt{.timer} en \texttt{.service}, puis \texttt{systemctl start <service>} par SSH\\
\textbf{Confirmation} & \textbf{Oui}, par la politique de l'application. Le docstring de l'outil dit pourtant « pas de confirmation requise » : l'outil se juge inoffensif, l'application le soumet quand même au clic humain\\
\textbf{Retour} & \texttt{\{"action": "backup\_started", "target", "status": "running"\}} ou \texttt{\{"error": ..., "action": "backup\_failed"\}}\\
\bottomrule
\end{longtable}

% ---------------------------------------------------
\subsection{\fn{create\_snapshot(volume\_id, volume\_group, size\_gb)}}

\begin{longtable}{@{}p{3.2cm}p{11.6cm}@{}}
\toprule
\endhead
\textbf{Où} & \file{app/mcp/storage/tools/snapshot.py}\\
\textbf{Qui} & Même circuit : proposition puis confirmation\\
\textbf{Paramètres} & \texttt{volume\_id} (obligatoire), \texttt{volume\_group} (déduit par \texttt{lvs} s'il est absent), \texttt{size\_gb=1.0}\\
\textbf{Effet réel} & Nom horodaté \texttt{<lv>-snap-AAAAMMJJ-HHMMSS}, puis \texttt{sudo lvcreate -{}-snapshot -{}-size <n>G -{}-name <nom> /dev/<vg>/<lv>}\\
\textbf{Confirmation} & Oui (appartenance à \texttt{REMEDIATION\_TOOLS}). L'outil lui-même n'exige pas \texttt{confirm=True}\\
\textbf{Nature} & Non destructive : « elle ajoute un volume, elle n'écrase rien »\\
\textbf{Garde-fou interne} & Si le logical volume n'existe pas : \texttt{\{"error": "Logical volume '<x>' not found on the VM"\}}, sans rien exécuter\\
\bottomrule
\end{longtable}

% ---------------------------------------------------
\subsection{\fn{restore\_from\_backup(target, backup\_id, confirm)} --- la plus dangereuse}

\begin{longtable}{@{}p{3.2cm}p{11.6cm}@{}}
\toprule
\endhead
\textbf{Où} & \file{app/mcp/storage/tools/restore.py}\\
\textbf{Classement} & Seul membre de \texttt{DESTRUCTIVE\_TOOLS}, et membre de \texttt{REMEDIATION\_TOOLS}\\
\textbf{Paramètres} & \texttt{target} (volume), \texttt{backup\_id} (nom du snapshot LVM), \texttt{confirm=False}\\
\textbf{Effet réel} & \texttt{sudo lvconvert -{}-merge <snapshot>} --- « elle écrase l'état courant du volume par celui du snapshot » et supprime le snapshot\\
\textbf{Double verrou} & \emph{(1)} interception par \texttt{REMEDIATION\_TOOLS} ; \emph{(2)} l'outil lui-même refuse d'agir sans \texttt{confirm=True}\\
\textbf{Sans confirmation} & Retourne \texttt{\{"action": "restore\_requires\_confirmation", "warning": "\ldots{} écrasera les données actuelles du volume et supprimera le snapshot"\}} --- \textbf{rien n'est exécuté}\\
\textbf{Note technique} & LVM ne fusionne qu'au prochain démontage/remontage : le retour le dit et invite à vérifier avec \texttt{lvs}\\
\bottomrule
\end{longtable}

\begin{lstlisting}[language=Python,caption={Le refus par défaut, dans l'outil lui-même}]
def restore_from_backup(target: str, backup_id: str, confirm: bool = False) -> dict:
    if not confirm:
        return {
            "action": "restore_requires_confirmation",
            "target": target, "backup_id": backup_id,
            "warning": (f"Fusionner le snapshot '{backup_id}' dans '{target}' ecrasera "
                        "les donnees actuelles du volume et supprimera le snapshot. "
                        "Confirmation explicite requise."),
        }
    output = run(f"sudo lvconvert --merge {quote(backup_id)}")
\end{lstlisting}

% ---------------------------------------------------
\subsection{\fn{repair\_backup\_repository(target, confirm)}}

C'est l'outil le plus élaboré du projet, et le seul qui \emph{modifie une
configuration système}. Il mérite un développement.

\subsubsection*{Le problème qu'il traite}
Le docstring l'expose : « Jusqu'ici OASIS savait relancer un job en échec, pas
corriger ce qui le faisait échouer. Sur un job dont le dépôt restic est
introuvable, relancer ne pouvait que reproduire la même erreur : l'ingénieur
voyait un diagnostic juste suivi d'une remédiation vouée à rater. »

\subsubsection*{Les trois principes déclarés}
\begin{enumerate}
  \item \textbf{Il ne devine rien.} Le dépôt de remplacement est \emph{lu sur la
  machine}, pas fourni par le modèle, et seuls les jobs dont la dernière
  exécution a réussi servent de référence.
  \item \textbf{Il ne touche pas au fichier d'unité.} Un drop-in s'annule d'un
  \texttt{systemctl revert}, ce que la réécriture d'une unité ne permet pas.
  \item \textbf{Il refuse plutôt que d'approximer.}
\end{enumerate}

\subsubsection*{Déroulé, étape par étape}

\begin{longtable}{@{}p{0.7cm}p{4.6cm}p{9.4cm}@{}}
\toprule
\textbf{\#} & \textbf{Étape} & \textbf{Détail}\\
\midrule
\endhead
1 & Lecture de la configuration & \texttt{systemctl show <service> -p Environment -p EnvironmentFiles -p FragmentPath}\\
2 & Unit inexistante & \texttt{\{"action": "repair\_failed", "error": "No systemd unit named \ldots"\}}\\
3 & Pas de \texttt{RESTIC\_REPOSITORY} direct & \texttt{"no\_repair\_possible"} : la configuration vient d'un fichier que le compte SSH ne peut pas lire\\
4 & Test du dépôt & \fn{\_repository\_state} lance \texttt{restic cat config} avec un \textbf{faux mot de passe} volontaire\\
5 & Interprétation & \texttt{"version"} dans la sortie \arrowto{} \texttt{valid} ; « wrong password » \arrowto{} \texttt{protected} ; « no such file » \arrowto{} \texttt{missing} ; sinon \texttt{unknown}\\
6 & Dépôt valide ou protégé & \texttt{"no\_repair\_needed"} : « The job fails for another reason --- read its journal with \fn{get\_backup\_logs} »\\
7 & Recherche du donneur & \fn{\_healthy\_donor} liste les services de sauvegarde, écarte les units de la distribution, lit leur \texttt{EnvironmentFiles}\\
8 & Critères stricts & Le fichier doit être celui du plus grand nombre de jobs, \textbf{strictement} (égalité \arrowto{} abstention), et au moins un job l'utilisant doit avoir réussi\\
9 & Pas de donneur & \texttt{"no\_repair\_possible"} avec la liste des voisins examinés\\
10 & \textbf{Sans \texttt{confirm}} & \texttt{"repair\_requires\_confirmation"} avec un champ \texttt{plan} en langage clair\\
11 & Avec \texttt{confirm=True} & Écriture du drop-in, \texttt{systemctl daemon-reload}, puis relecture pour vérifier la prise en compte\\
12 & Refus de sudo & \fn{\_missing\_privilege} traduit l'erreur en consigne exploitable (règle sudoers précise, ou marche à suivre manuelle)\\
\bottomrule
\end{longtable}

\begin{keybox}
\textbf{L'astuce du faux mot de passe.} \fn{\_repository\_state} passe
\texttt{-{}-password-command 'printf x'} exprès : « restic lit le fichier
\texttt{config} du dépôt avant de vérifier la clé, donc l'absence de dépôt et le
mauvais mot de passe produisent deux erreurs distinctes. C'est ce qui permet de
distinguer "rien ici" de "dépôt bien présent" sans manipuler le moindre
secret. » C'est un excellent exemple de conception défensive à citer.
\end{keybox}

\subsubsection*{Réversibilité}
Le drop-in écrit contient un en-tête qui identifie son auteur et donne la
commande d'annulation :

\begin{lstlisting}[caption={Contenu du fichier /etc/systemd/system/<service>.d/50-oasis-repair.conf}]
# Written by OASIS on 2026-08-23 14:12 UTC (repair_backup_repository).
# backup-app-critical.service pointed at /srv/repo, which is not a restic repository.
# It now reads the environment shared by this host's healthy backup jobs.
# Undo with: sudo systemctl revert backup-app-critical.service
[Service]
EnvironmentFile=/etc/oasis/restic.env
\end{lstlisting}

% =====================================================================
\section{Le mécanisme de confirmation humaine}

\subsection{Où il est implémenté}

\begin{longtable}{@{}p{4.4cm}p{10.4cm}@{}}
\toprule
\textbf{Couche} & \textbf{Implémentation}\\
\midrule
\endhead
Classification & \texttt{REMEDIATION\_TOOLS} dans \file{mcp/storage/server.py}\\
Interception & \fn{StoragePersona.\_propose\_action} dans \file{personas/storage/agent.py}\\
Persistance & Table \tbl{action\_requests}, statut \texttt{pending}\\
Rattachement & \tbl{action\_requests.chat\_message\_id} via \fn{chat\_service.attach\_actions}\\
Exposition & \cls{ActionRequestOut} dans \cls{ChatMessageOut} et \cls{IncidentTicketDetail}\\
Interface & \cls{ActionProposal} (\file{frontend/src/components/ActionProposal.tsx})\\
Exécution & \route{POST /storage/actions/\{id\}/confirm} et \fn{confirm\_and\_execute\_action}\\
Refus & \route{POST /storage/actions/\{id\}/reject} + note d'incident\\
Vérification & \fn{\_close\_the\_loop} + \file{services/incident\_verification.py}\\
\bottomrule
\end{longtable}

\subsection{Le raisonnement derrière la séparation \texttt{DESTRUCTIVE} / \texttt{REMEDIATION}}

Le commentaire de \file{server.py} explique pourquoi deux ensembles et non un :

\begin{quote}
\itshape
« La distinction avec \texttt{DESTRUCTIVE\_TOOLS} n'est pas la même : un restore
est destructif, relancer une sauvegarde ne l'est pas --- mais les deux changent
l'état du système, et aucun des deux ne doit partir sur la seule décision du
modèle. \texttt{DESTRUCTIVE\_TOOLS} continue de qualifier le risque dans
l'audit. »
\end{quote}

\noindent
Autrement dit : \texttt{REMEDIATION\_TOOLS} répond à « faut-il un humain ? »
(oui pour les quatre), \texttt{DESTRUCTIVE\_TOOLS} répond à « quel niveau de
risque ? » (HIGH pour un seul).

\subsection{Ce qui rend le contournement difficile}

\begin{enumerate}
  \item \textbf{Le modèle ne peut pas exécuter directement.} L'aiguillage est un
  test d'appartenance à un ensemble figé, en dur dans le code.
  \item \textbf{Écrire l'appel en texte ne sert à rien.} Un JSON recopié dans la
  réponse ne crée rien --- et \fn{\_recover\_written\_call} ne le transforme
  qu'en \emph{proposition}, avec double validation de l'outil et de la cible.
  \item \textbf{Le prompt le dit au modèle.} « never say a remediation has been
  launched: it is waiting on a human decision. »
  \item \textbf{Le résultat renvoyé au modèle le répète.} « Action NOT executed
  \ldots{} Present it as a recommendation, never as a fact. »
  \item \textbf{Sans utilisateur identifié, refus net.} Si \texttt{created\_by}
  est \texttt{None}, l'agent renvoie
  \texttt{\{"status": "denied", "reason": "confirmation\_required"\}} --- « on
  refuse plutôt que de l'attribuer à quelqu'un au hasard ».
  \item \textbf{Double exécution impossible.} Le test
  \texttt{status != PENDING} lève une \cls{ValueError}, traduite en 404.
  \item \textbf{Deux outils exigent en plus \texttt{confirm=True}} au niveau de
  l'implémentation.
\end{enumerate}

\begin{warnbox}
\textbf{La limite honnête à énoncer.} Les routes de \file{actions.py} n'exigent
que \fn{get\_current\_user}, pas \fn{require\_role}. Un utilisateur authentifié
avec le rôle \texttt{linux\_engineer} ne peut ni ouvrir le chat ni voir les
incidents (ces routes sont protégées), mais s'il obtenait un identifiant
d'action valide, la route \route{/confirm} ne le rejetterait pas. Ajouter
\fn{require\_role} sur ce routeur serait le durcissement le plus rentable du
projet.
\end{warnbox}

\subsection{Le workflow complet, en une table}

\begin{longtable}{@{}p{0.7cm}p{3.7cm}p{4.3cm}p{5.7cm}@{}}
\toprule
\textbf{\#} & \textbf{Étape} & \textbf{Acteur} & \textbf{Trace laissée}\\
\midrule
\endhead
1 & Demande utilisateur & Ingénieur & \tbl{chat\_messages} (rôle \texttt{user})\\
2 & Lecture de la VM & Agent + MCP + SSH & \tbl{mcp\_calls} (\texttt{success})\\
3 & Analyse & Groq & \tbl{chat\_messages} (rôle \texttt{assistant})\\
4 & Proposition & \fn{\_propose\_action} & \tbl{action\_requests} (\texttt{pending})\\
5 & Rattachement & \fn{attach\_actions} & \texttt{chat\_message\_id} renseigné\\
6 & Affichage & \cls{ActionProposal} & ---\\
7 & \textbf{Décision humaine} & \textbf{Ingénieur} & \texttt{confirmed\_by}, \texttt{confirmed\_at}\\
8 & Exécution & MCP + SSH & \tbl{mcp\_calls} avec \texttt{action\_request\_id}\\
9 & Résultat & \fn{confirm\_and\_execute\_action} & \texttt{status}, \texttt{result} (JSONB), \texttt{completed\_at}\\
10 & Vérification & \fn{verify} & \tbl{mcp\_calls} (relecture) + \tbl{incident\_ticket\_notes}\\
11 & Clôture & \fn{resolve\_ticket} & \texttt{status='resolved'}, \texttt{resolved\_at}, note\\
\bottomrule
\end{longtable}

\begin{keybox}
\textbf{Onze étapes, sept traces en base.} C'est ce qui permet de répondre, six
mois plus tard, à la question « qui a lancé cette restauration, sur la foi de
quel relevé, et qu'est-ce que ça a donné ? ».
\end{keybox}
